%=============================================================================
% APPENDIX T: FAMILY AND CLOCK-TYPE PARAMETRIZATION
% DFD Unified Review v3.0
%=============================================================================

\section{Family and Clock-Type Parametrization of LPI Tests}\label{app:family-clock}

This appendix presents a phenomenological parametrization organizing 
clock comparison tests by chemical family and clock type. The framework 
provides a compact way to encode where current data pull and where 
future tests should focus.

%-----------------------------------------------------------------------------
\subsection{Two-Parameter Model}\label{app:fc-model}
%-----------------------------------------------------------------------------

Motivated by the pattern of hints and nulls in clock comparisons, we 
parameterize the gravitational coupling coefficient as:
\begin{equation}\label{eq:K-family-clock}
K_i = k_N C_N^{(i)} + k_e C_e^{(i)} + k_\alpha \kappa_\alpha^{(i)},
\end{equation}
where:
\begin{itemize}
\item $C_N^{(i)}$: Nuclear-sector charge depending on chemical family
\item $C_e^{(i)}$: Electronic-sector charge depending on clock type
\item $\kappa_\alpha^{(i)} = S_i^\alpha$: Standard $\alpha$-sensitivity
\item $k_N, k_e, k_\alpha$: Coupling strengths to be fit or constrained
\end{itemize}

\paragraph{Family charges.}
Based on chemical grouping:
\begin{center}
\begin{tabular}{lc}
\toprule
\textbf{Element family} & $C_N$ \\
\midrule
Alkaline earth (Sr, Ca, Mg) & 0 \\
Alkali (Cs, Rb, H) & 1 \\
Post-transition (Al, Hg, In) & 1.5 \\
Lanthanide (Yb, Dy) & 2 \\
\bottomrule
\end{tabular}
\end{center}

\paragraph{Clock-type charges.}
Based on interrogation mode:
\begin{equation}
C_e = \begin{cases}
0 & \text{optical neutral} \\
0.5 & \text{trapped ion} \\
1 & \text{microwave hyperfine}
\end{cases}
\end{equation}

These assignments are deliberately coarse; the point is not that nuclear 
scalar charges take precisely these values, but that grouping by family 
and clock type yields a testable pattern.

%-----------------------------------------------------------------------------
\subsection{Constraints from Data}\label{app:fc-constraints}
%-----------------------------------------------------------------------------

For each clock pair $(A, B)$, the observable is:
\begin{equation}
\Delta K_{AB} = k_N \Delta C_N + k_e \Delta C_e + k_\alpha \Delta\kappa_\alpha,
\end{equation}
where $\Delta C_N = C_N^{(A)} - C_N^{(B)}$ and similarly for other quantities.

\paragraph{E3/E2 constraint on $k_\alpha$.}
The Yb$^+$ E3/E2 same-ion comparison has $\Delta C_N = \Delta C_e = 0$ but 
$\Delta\kappa_\alpha = -6.95$. The PTB bound $|\Delta K_{\rm E3/E2}| < 10^{-8}$ 
thus constrains:
\begin{equation}
|k_\alpha| < 1.4 \times 10^{-9}.
\end{equation}
This effectively forces $k_\alpha \to 0$, eliminating pure-$\alpha$ 
coupling from the model.

\paragraph{Cross-species constraints.}
With $k_\alpha = 0$ fixed, the two-parameter model $(k_N, k_e)$ is 
constrained by:
\begin{itemize}
\item H/Cs null ($\Delta C_N = \Delta C_e = 0$): automatically satisfied
\item Hg$^+$/Cs: $\Delta C_N = 0.5$, $\Delta C_e = -0.5$, bound $|\Delta K| < 5.8 \times 10^{-6}$
\item Dy/Cs: $\Delta C_N = 1$, $\Delta C_e = -1$, bound $|\Delta K| < 10^{-5}$
\item Cs/Sr hint: $\Delta C_N = 1$, $\Delta C_e = 1$, suggests $\Delta K \sim 3 \times 10^{-5}$
\end{itemize}

A joint fit yields $k_N \sim 6 \times 10^{-6}$, $k_e \sim 1.5 \times 10^{-5}$ 
with $\chi^2_\nu \approx 1$.

%-----------------------------------------------------------------------------
\subsection{Predictions for Untested Channels}\label{app:fc-predictions}
%-----------------------------------------------------------------------------

The model predicts specific $\Delta K$ values for channels where 
high-precision ratios exist but LPI analyses have not been performed:

\begin{table}[h!]
\centering
\caption{Family+clock model predictions for untested LPI channels.}\label{tab:fc-predictions}
\footnotesize
\renewcommand{\arraystretch}{1.2}
\begin{tabular}{lcccc}
\toprule
\textbf{Channel} & $\Delta C_N$ & $\Delta C_e$ & \textbf{Predicted} $\Delta K$ & \textbf{Test type} \\
\midrule
Sr$^+$/Sr & 0 & $+0.5$ & $7.5 \times 10^{-6}$ & Pure electronic \\
Ca$^+$/Sr & 0 & $+0.5$ & $7.5 \times 10^{-6}$ & Pure electronic \\
Hg/Sr & $+1.5$ & 0 & $9 \times 10^{-6}$ & Pure nuclear-family \\
Yb$^+$/Sr$^+$ & $+2$ & 0 & $1.2 \times 10^{-5}$ & Pure nuclear-family \\
Hg/Yb & $-0.5$ & 0 & $-3 \times 10^{-6}$ & Partial cancellation \\
Ca/Sr & 0 & 0 & 0 & Null prediction \\
\bottomrule
\end{tabular}
\end{table}

\paragraph{Falsification criteria.}
\begin{enumerate}
\item An observed Ca/Sr LPI signal at $\sim 10^{-5}$ would falsify the family structure.
\item Hg/Sr or Ca$^+$/Sr showing null results at $< 10^{-6}$ would severely constrain both $k_N$ and $k_e$.
\item Consistency across untested channels validates the two-parameter structure.
\end{enumerate}

%-----------------------------------------------------------------------------
\subsection{Relation to DFD Microsector}\label{app:fc-dfd}
%-----------------------------------------------------------------------------

The phenomenological charges $(C_N, C_e)$ are not derived from first 
principles in this appendix. However, they are compatible with the 
DFD microsector in the following sense:
\begin{itemize}
\item The overall scale $k_N, k_e \sim 10^{-5}$--$10^{-6}$ is consistent 
with $k_\alpha = \alpha^2/(2\pi)$ structure.
\item The family grouping (alkaline earth vs.\ lanthanide) suggests 
coupling to properties correlated with atomic structure, not just $\alpha$.
\item The clock-type structure (ion vs.\ neutral) aligns with the 
sector-coupling hierarchy in Sec.~\ref{sec:clocks-hierarchy}.
\end{itemize}

A full derivation of $(C_N, C_e)$ from the $\mathbb{CP}^2 \times S^3$ microsector 
remains an open problem. The present appendix establishes the empirical 
pattern that such a derivation must reproduce.

%-----------------------------------------------------------------------------
\subsection{Summary}\label{app:fc-summary}
%-----------------------------------------------------------------------------

The family+clock parametrization provides:
\begin{enumerate}
\item A compact organization of existing LPI constraints
\item Specific predictions for channels where analyses are actionable
\item Clear falsification criteria
\item A target pattern for microsector derivation
\end{enumerate}

The decisive tests are Hg/Sr (pure nuclear-family) and Sr$^+$/Sr 
(pure electronic), both of which can be performed with existing clock 
technology.
